\paragraph{IV 21, \S2: relative divisors on projective and multiprojective bundles (cote25, sites 106--111).}
The section applies the classification of invertible Modules from1.1 to positive relative divisors. At sites 106--107 the canonical direct-image formula marked$(*)$ converts the section defining a divisor into a map involving a symmetric power of $\underline E$; the author expresses fiberwise regularity through universal injectivity and the surjectivity of its transpose. Proposition 2.2 describes the resulting correspondence with invertible quotients of the dual symmetric power, and Corollaries2.3--2.4 give the representing projective bundles and their sum over the degree. Paragraph2.5 introduces a finite family of projective bundles and their fiber product. Proposition 2.6 and paragraph 2.7 extend the classification and degree notation to this multiprojective bundle, explicitly treating the partial degrees as locally constant functions on the base. Sites108--109 connect this construction to the planned paragraph 25 on the resultant and discriminant of forms. Paragraph2.8 invokes ChapterII, paragraph 2, in the displayed direct-image formula$(xx)$; its canceled alternative reference to Künneth is readable. Proposition 2.9 and Corollary 2.10 describe divisors of a fixed multidegree by invertible quotients of a tensor product of dual symmetric powers and give the associated representing projective bundle and sum. Paragraph2.11 then extracts the relative representability argument from the projective-bundle case and relates it to the Picard prescheme. At site 110, under the displayed hypothesis on $f_*(\underline O_X)$, the proof uses ChapterIII, paragraph 7, to introduce the finitely presented Module $\underline Q$ representing the relevant functor, identifies sections with homomorphisms to an invertible Module, and states the correspondence with sections of $\underline P(\underline Q)$. The proposition number on that page is typed2.11 with an uncertain pen correction resembling2 on its last digit; site 111 explicitly refers to2.12. Theorem 2.13 at site 111 gives the open retrocompact representing sub-prescheme, and the text identifies the entire projective bundle in its geometrically integral-fiber case. Its proof invokes the openness of transverse regularity through the incomplete reference ``par.11 ...'' and obtains retrocompactness by reduction to the noetherian case. The transcription retains several visible features that affect reading: site 106 says ``section de $\underline E$'' and uses $\underline O_X(\underline D)$; site 107 prints ``décomposition de D'' and has a crowded extra subscript and crossed group in the symmetric-power expression of2.4; the displayed isomorphism of2.6 at site 108 omits the $\underline M$ factor mentioned in the surrounding statement; and site 109 prints $\underline O_X(D)$ and $f_*(\underline D)$ in the two formulas. These are preserved with the scanned alphabets and indices. The N/G variable at site 110, the corrected proposition digit, and several typewriter overprints remain explicitly uncertain.
